%% TRANSCRIBED (PDF 146-153)
%% NOTE: the tail of 1.4 (end of the proof of the theorem and its Remark,
%% printed on IV-9 = PDF 154 above the §2 heading) is transcribed at the
%% end of this file so that §1 is complete; ch4_s2.tex starts at §2.

%% ---------- Chapter IV opening (printed page IV-1) ----------
Let $K$ be a number field and let $E$ be an elliptic curve over $K$.
If $\ell$ is a prime number, let
\[
\rho_\ell : \Gal(\bar K/K) \longrightarrow \Aut(V_\ell(E))
\]
be the corresponding \textit{l}-adic representation of $K$,
cf.~chap.~\ref{chap:I}, \S\ref{sec:I.1.2}.

The main result of this Chapter is the determination of the Lie algebra
of the \textit{l}-adic Lie group $G_\ell = \Ima(\rho_\ell)$.  This is
based on a finiteness theorem of \v{S}afarevi\v{c} (\S\ref{sec:IV.1.4})
combined with the properties of locally algebraic abelian
representations (chap.~\ref{chap:III}) and Tate's local theory of
elliptic curves with non-integral modular invariant
(Appendix, \S\ref{sec:IV.A.1}).  The variation of $G_\ell$ with $\ell$
is studied in \S\ref{sec:IV.3}.

The Appendix gives analogous results in the local case (i.e.\ when $K$
is a local field).

\sechead{1}{Preliminaries}{PRELIMINARIES}

\subsechead{1.1}{Elliptic curves
  \texorpdfstring{(cf.~Cassels~\cite{cassels-dio}, Deuring~\cite{deuring},
  Igusa~\cite{igusa})}{}}
%% NOTE: the printed heading (IV-2) carries the trailing note
%% ``(cf. Cassels [5], Deuring [9], Igusa [10])''; the book's own table of
%% contents lists only ``1.1  Elliptic curves''.  The note is therefore
%% kept in the body heading (with live \cite's), stripped from the PDF
%% bookmark; it will appear in the document ToC line.

By an elliptic curve, we mean an abelian variety of dimension $1$,
i.e. a complete, non singular, connected curve of genus $1$ with a
given rational point $P_o$, taken as an origin for the composition law
(and often written $0$).
%% NOTE: the book prints the subscript of $P$ as a small round lowercase
%% ``o'' (typewriter nought), clearly distinct from the oval digit ``0''
%% of the same sentence; transcribed literally as $P_o$ (cf. the $K_o$
%% note in ch1_s1.tex).

Let $E$ be such a curve.  It is well known that $E$ may be embedded, as
a non-singular cubic, in the projective plane $\mathbf{P}_{2/K}$, in
such a way that $P_o$ becomes a ``flex'' (one takes the projective
embedding defined by the complete linear series containing the divisor
$3.P_o$).  In this embedding, three points $P_1$, $P_2$, $P_3$ have sum
$0$ if and only if the divisor $P_1 + P_2 + P_3$ is the intersection of
$E$ with a line.  By choosing a suitable coordinate system, the
equation of $E$ can be written in Weierstrass form
\[
y^2 = 4x^3 - g_2x - g_3
\]
where $x$, $y$ are non-homogeneous coordinates and the origin $P_o$ is
the point at infinity on the $y$-axis.  The discriminant
\[
\Delta = g_2^3 - 27g_3^2
\]
is non-zero.

The coefficients $g_2$, $g_3$ are determined up to the transformations
$g_2 \longmapsto u^4g_2$, $g_3 \longmapsto u^6g_3$, $u \in K^*$.  The
modular invariant $j$ of $E$ is
\[
j = 2^63^3\,\frac{g_2^3}{g_2^3 - 27g_3^2}
  = 2^63^3\,\frac{g_2^3}{\Delta}
\]

Two elliptic curves have the same $j$ invariant if and only if they
become isomorphic over the algebraic closure of $K$.

(All this remains valid over an arbitrary field, except that, when the
characteristic is $2$ or $3$, the equation of $E$ has to be written in
the more general form
\[
y^2 + a_1xy + a_3y + x^3 + a_2x^2 + a_4x + a_6 = 0\,.
\]
Here again, $0$ is the point at infinity on the $y$-axis and the
corresponding tangent is the line at infinity.  There are corresponding
definitions for $\Delta$ and $j$, for which we refer to
Deuring~\cite{deuring} or Ogg~\cite{ogg}; note, however, that there is
a misprint in Ogg's formula for $\Delta$: the coefficient of
$\beta_4^3$ should be $-8$ instead of $-1$.)

\subsechead{1.2}{Good reduction}

Let $v \in \Sigma_K$ be a place of the number field $K$.  We denote by
$O_v$ (resp.\ $\underline{m}_v$, $k_v$) the corresponding local ring in
$K$ (resp.\ its maximal ideal, its residue field).
%% NOTE: the book prints the maximal ideal as an underlined italic ``m''
%% (the underline standing for the boldface ideal symbol); transcribed
%% literally as $\underline{m}_v$.

Let $E$ be an elliptic curve over $K$.  One says that $E$ has
\emph{good reduction at} $v$ if one can find a coordinate system in
$\mathbf{P}_{2/K}$ such that the corresponding equation $f$ for $E$ has
coefficient in $O_v$ and its reduction $\tilde{f} \bmod
\underline{m}_v$ defines a \emph{non-singular} cubic
$\widetilde{E}_v$ (hence an elliptic curve) over the residue field
$k_v$ (in other words, the discriminant $\Delta(f)$ of $f$ must be an
invertible element of $O_v$).  The curve $\widetilde{E}_v$ is called
the \emph{reduction} of $E$ at $v$; it does not depend on the choice of
$f$, provided, of course, that $\Delta(f) \in O_v^*$.

One can prove that the above definition is equivalent to the following
one: there is an abelian scheme $E_v$ over $\Spec(O_v)$, in the sense
of Mumford~\cite{mumford-git}, chap.~VI, whose generic fiber is $E$;
this scheme is then unique, and its special fiber is
$\widetilde{E}_v$.  Note that $\widetilde{E}_v$ is defined over the
finite field $k_v$; we denote its \emph{Frobenius endomorphism} by
$F_v$.

On either definition, one sees that $E$ has \emph{good reduction for
almost all places of} $K$.

If $E$ has good reduction at a given place $v$, its $j$ invariant is
\emph{integral at} $v$ (i.e. belongs to $O_v$) and its reduction
$\tilde{j} \bmod \underline{m}_v$ is the $j$ invariant of the reduced
curve $\widetilde{E}_v$.

The converse is almost true, but not quite: if $j$ belongs to $O_v$,
there is a finite extension $L$ of $K$ such that $E \times_K L$ has
good reduction at all the places of $L$ dividing $v$ (this is the
``potential good reduction'' of Serre-Tate~\cite{serre-tate-good},
\S2).  For the proof of this, see Deuring~\cite{serre-corps}, \S4,
n°~3.
%% NOTE: the book literally prints ``Deuring [29]'' here, although [29]
%% is Serre's Corps Locaux and Deuring is [9] in the bibliography (cited
%% in 1.1 above).  This is a misprint in the original; it is kept as
%% printed, and \cite{serre-corps} is the key that prints [29].

\runhead{Remark}

The definitions and results of this section have nothing to do with
number fields.  They apply to every field with a discrete valuation.

\subsechead{1.3}{Properties of $V_\ell$ related to good reduction}
%% NOTE: the skeleton heading had $V_l$; corrected to $V_\ell$ as printed
%% (IV-4; the skeleton was generated from the ToC dump).

Let $\ell$ be a prime number.  We define, as in chap.~\ref{chap:I},
\S\ref{sec:I.1.2}, the Galois modules $T_\ell$ and $V_\ell$ by:
\[
V_\ell = T_\ell \tens \Q_\ell, \quad T_\ell = \varprojlim\ E_{\ell^n}
\]
where $E_{\ell^n}$ is the kernel of
$\ell^n : E(\bar K) \longrightarrow E(\bar K)$.

We denote by $\rho_\ell$ the corresponding homomorphism of
$\Gal(\bar K/K)$ into $\Aut(T_\ell)$.  Recall that $E_{\ell^n}$,
$T_\ell$ and $V_\ell$ are of rank $2$ over $\Z/\ell^n\Z$, $\Z_\ell$ and
$\Q_\ell$, respectively.

Let now $v$ be a place of $K$, with $p_v \neq \ell$ and let $\bar v$ be
some extension of $v$ to $\bar K$; let $D$ (resp.\ $I$) be the
corresponding decomposition group (resp.\ inertia group),
cf.~chap.~\ref{chap:I}, \S\ref{sec:I.2.1}.  If $E$ has good reduction
at $v$, one easily sees that reduction at $\bar v$ defines an
\emph{isomorphism} of $E_{\ell^n}$ onto the corresponding module for
the reduced curve $\widetilde{E}_v$.  In particular, $E_{\ell^n}$,
$T_\ell$, $V_\ell$ are \emph{unramified at} $v$
(chap.~\ref{chap:I}, \S\ref{sec:I.2.1}) and the Frobenius automorphism
$F_{v,\rho_\ell}$ of $T_\ell$ corresponds to the Frobenius endomorphism
$F_v$ of $\widetilde{E}_v$.  Hence:
\[
\det(F_{v,\rho_\ell}) = \det(F_v) = N_v
\]
and
\[
\det(1-F_{v,\rho_\ell}) = \det(1-F_v) = 1 - \mathrm{Tr}(F_v) + N_v
\]
is equal to the number of $k_v$-points of $\widetilde{E}_v$.

Conversely:

\medskip
\noindent\textbf{CRITERION OF N\'ERON-OGG-\v{S}AFAREVI\v{C}.}%
\hspace{0.6em}\emph{If} $V_\ell$ is \emph{unramified at} $v$ \emph{for
some} $\ell \neq p_v$, \emph{then} $E$ \emph{has good reduction at}
$v$.
\medskip
%% NOTE: the book prints the criterion as a bold run-in heading with a
%% roman statement body; the original emphasis (underline = italic) is
%% transcribed with \emph.  No statement number is printed, and the
%% proof remark below refers to it simply as ``the theorem''.

For the proof, see Serre-Tate~\cite{serre-tate-good}, \S1.

\begin{corollary}
Let $E$ and $E'$ be two elliptic curves which are isogenous (over $K$).
If one of them has good reduction at a place $v$, the same is true for
the other one.
\end{corollary}

(Recall that $E$ and $E'$ are said to be \emph{isogenous} if there
exists a non-trivial morphism $E \longrightarrow E'$.)

This follows from the theorem, since the \textit{l}-adic
representations associated with $E$ and $E'$ are isomorphic.
%% NOTE: ``the theorem'' = the Criterion of Néron-Ogg-Šafarevič above
%% (unnumbered in the book).

\runhead{Remark}

For a direct proof of this corollary, see
Koizumi-Shimura~\cite{koizumi-shimura}.

\runhead{Exercise}

Let $S$ be the finite set of places where $E$ does not have good
reduction.  If $v \in \Sigma_K - S$, we denote by $t_v$ the number of
$k_v$-points of the reduced curve $\widetilde{E}_v$.

\begin{alist}
\item Let $\ell$ be a prime number and let $m$ be a positive integer.
Show that the following properties are equivalent:
\begin{rlist}
\item $t_v \equiv 0 \bmod \ell^m$ for all $v \in \Sigma_K - S$,
$p_v \neq \ell$.
\item The set of $v \in \Sigma_K - S$ such that $t_v \equiv 0 \bmod
\ell^m$ has density one (cf.~chap.~\ref{chap:I},
\S\ref{sec:I.2.2}).
\item For all $s \in \Ima(\rho_\ell)$, one has $\det(1-s) \equiv 0
\bmod \ell^m$.
\end{rlist}
(The equivalence of (ii) and (iii) follows from \v{C}ebotarev's
density theorem.  The implications (i) $\Rightarrow$ (ii) and (iii)
$\Rightarrow$ (i) are easy.)

\item We take now $m = 1$.  Show that the properties (i), (ii), (iii)
are equivalent to:

(iv) There exists an elliptic curve $E'$ over $K$ such that:
%% NOTE: the book labels the two items below with greek letters $(\alpha)$,
%% $(\beta)$ -- not (a), (b) as the OCR dump suggested.
\begin{enumerate}[leftmargin=2.2em,itemsep=0.15ex,topsep=0.4ex]
\item[$(\alpha)$] Either $E'$ is isomorphic to $E$, or there exist an
isogeny $E' \longrightarrow E$ of degree $\ell$.
\item[$(\beta)$] The group $E'(K)$ contains an element of order
$\ell$.
\end{enumerate}
(The implication (iv) $\Rightarrow$ (iii) is easy.  For the proof of
the converse, use Exer.~2 of chap.~\ref{chap:I},
\S\ref{sec:I.1.1}.) $\longrightarrow$ [for $m \geq 2$, see
Katz~\cite{katz-torsion}.]
\end{alist}

\subsechead{1.4}{\v{S}afarevi\v{c}'s theorem}

It is the following (cf.~\cite{safarevic}):

\begin{theorem}
Let $S$ be a finite set of places of $K$.  The set of isomorphism
classes of elliptic curves over $K$, with good reduction at all places
not in $S$, is finite.
\end{theorem}

Since isogenous curves have the same bad reduction set
(cf.~\S\ref{sec:IV.1.3}), this implies:

\begin{corollary}
Let $E$ be an elliptic curve over $K$.  Then, up to isomorphism, there
are only a finite number of elliptic curves which are $K$-isogenous to
$E$.
\end{corollary}

To prove the theorem, we use the following criterion for good
reduction:

\begin{lemma}
Let $S$ be a finite set of places of $K$ containing the divisors of
$2$ and $3$, and such that the ring $O_S$ of $S$-integers is
principal.  Then, an elliptic curve $E$ defined over $K$ has good
reduction outside $S$ if and only if its equation can be put in the
Weierstrass form $y^2 = 4x^3 - g_2x - g_3$ with $g_i \in O_S$ and
$\Delta = g_2^3 - 27g_3^2 \in O_S^*$ (the group of units of $O_S$).
\end{lemma}

\begin{proof}
The sufficiency is trivial.  To prove necessity, we write the curve
$E$ in the form
\begin{equation}\tag{$*$}\label{eq:IV.1.4.star}
y^2 = 4x^3 - g'_2x - g'_3
\end{equation}
with $g'_i \in K$.  Let $v$ be a place of $K$ not in $S$.  Then, since
there is good reduction at $v$, and since the divisors of $2$ and $3$
do not belong to $S$, the curve $E$ can be written in the form
\[
y^2 = 4x^3 - g_{2,v}x - g_{3,v}
\]
with $g_{i,v}$ in the local ring at $v$ and the discriminant
$\Delta_v$ a unit in this ring.  Using the properties of the
Weierstrass form, there is an element $u_v \in K^*$ such that
$g_{2,v} = u_v^4g'_2$, $g_{3,v} = u_v^6g'_3$, $\Delta_v =
u^{12}\Delta'$;
%% NOTE: in $\Delta_v = u^{12}\Delta'$ the book prints $u$ without the
%% subscript $v$ (unlike $u_v^4$, $u_v^6$ just before; checked at 600 dpi
%% on IV-8); kept as printed.
moreover, as we can take $g_{i,v} = g'_i$ for almost all $v$, we see
that we can assume that $u_v = 1$ for almost all $v \notin S$.  Since
the ring $O_S$ is principal, there is an element $u \in K^*$ with
$v(u) = v(u_v)$ for all $v \notin S$.  Then, if we replace $x$ by
$u^{-2}x$ and $y$ by $u^{-3}y$ in (\ref{eq:IV.1.4.star}), the curve
$E$ takes the form
\[
y^2 = 4x^3 - g_2x - g_3
\]
with $g_2 = u^4g'_2$, $g_3 = u^6g'_3$ and $\Delta = u^{12}\Delta'$.
Since, by construction, $g_i \in O_S$ and $\Delta \in O_S^*$, the
lemma is established.
\end{proof}

\begin{proof}[Proof of the theorem.]
After possibly adding a finite number of places of $K$ to $S$, we may
assume that $S$ contains all the divisors of $2$ and $3$, and that the
ring $O_S$ is principal.  If $E$ is an elliptic curve defined over $K$
having good reduction outside $S$, the above lemma tells us that we
can write $E$ in the form
\begin{equation}\tag{$*$}\label{eq:IV.1.4.starB}
y^2 = 4x^3 - g_2x - g_3
\end{equation}
with $g_i \in O_S$ and $\Delta = g_2^3 - 27g_3^2 \in O_S^*$.  But,
since we are free to multiply $\Delta$ by any $u \in (O_S^*)^{12}$,
and since $O_S^*/(O_S^*)^{12}$ is a finite group, we see that there is
a finite set $X \subset O_S^*$ such that any elliptic curve of the
above type can be written in the form (\ref{eq:IV.1.4.starB}) with
$g_i \in O_S$ and $\Delta \in X$.  But, for a given $\Delta$, the
equation
\[
U^3 - 27V^2 = \Delta
\]
represents an affine elliptic curve.  Using a theorem of Siegel
(generalized by Mahler and Lang, cf.~Lang~\cite{lang-dio}, chap.~VII),
one sees that this equation has only a \emph{finite} number of
solutions in $O_S$.  This finishes the proof of the theorem.
\end{proof}
%% NOTE: the last two displays above are printed on IV-9 = PDF 154, which
%% continues §1.4 before the §2 heading; transcribed here so that the
%% section ends complete.  The pointer ``$\longrightarrow$ [for $m \geq 2$,
%% see Katz [64].]'' of the Exercise is the book's arrow convention.

\runhead{Remark}

There are many ways in which one can deduce \v{S}afarevi\v{c}'s
theorem from Siegel's.  The one we followed has been shown to us by
Tate.
